% !TEX program = pdflatex
% English CV for Ruocheng Wang.
% Adapted from the Sourabh Bajaj resume template supplied by the author.

\documentclass[letterpaper,11pt]{article}
\usepackage[T1]{fontenc}
\usepackage{lmodern}
\usepackage[margin=0.60in,top=0.52in,bottom=0.60in]{geometry}
\usepackage{microtype}
\usepackage{titlesec}
\usepackage{enumitem}
\usepackage[hidelinks]{hyperref}
\usepackage{fancyhdr}
\usepackage{tabularx}
\usepackage{xcolor}

\hypersetup{
  pdftitle={Ruocheng Wang - Curriculum Vitae},
  pdfauthor={Ruocheng Wang},
  pdfsubject={Academic curriculum vitae}
}
\pagestyle{fancy}
\fancyhf{}
\renewcommand{\headrulewidth}{0pt}
\renewcommand{\footrulewidth}{0pt}
\fancyfoot[L]{\footnotesize\color{gray}Ruocheng Wang}
\fancyfoot[R]{\footnotesize\color{gray}\thepage}
\setlength{\footskip}{23pt}
\setlength{\parindent}{0pt}
\setlength{\parskip}{0pt}
\setlength{\tabcolsep}{0pt}
\setlength{\emergencystretch}{1.5em}
\urlstyle{same}
\raggedbottom
\raggedright

\titleformat{\section}{\large\bfseries}{}{0pt}{}[\vspace{-4pt}\titlerule]
\titlespacing*{\section}{0pt}{10pt}{6pt}
\setlist[itemize]{leftmargin=1.2em,label=\textbullet,itemsep=3pt,topsep=3pt,parsep=0pt,partopsep=0pt}

\newcommand{\education}[3]{%
  \begin{tabularx}{\textwidth}{@{}Xr@{}}
    \textbf{#1} & #2 \\
  \end{tabularx}\par
  #3\par\vspace{5pt}
}
\newcommand{\publication}[3]{%
  \begingroup
  \setlength{\parskip}{0pt}
  \textbf{#1}\par
  #2\par
  \textit{#3}\par
  \endgroup\vspace{5pt}
}
\newcommand{\project}[2]{%
  \begin{tabularx}{\textwidth}{@{}>{\raggedright\arraybackslash}X@{\hspace{1.2em}}r@{}}
    \textbf{#1} & #2
  \end{tabularx}\par\vspace{1pt}
}

\begin{document}
\small

\begin{tabularx}{\textwidth}{@{}Xr@{}}
  \href{https://wang-ruocheng.github.io/}{\LARGE\textbf{Ruocheng Wang}} &
  \href{mailto:ignite@sjtu.edu.cn}{ignite@sjtu.edu.cn} \\
  \href{https://wang-ruocheng.github.io/}{wang-ruocheng.github.io} &
  +86 138-5136-2201
\end{tabularx}

\section{Education}
\education{Shanghai Jiao Tong University \& Shanghai Innovation Institute}
  {Sep. 2025 -- Present}
  {Direct-entry Ph.D. program (joint training).\\
   Advisor: Prof. Junchi Yan, School of Artificial Intelligence.}
\education{Shanghai Jiao Tong University}
  {Sep. 2019 -- Jun. 2025}
  {B.Sc. in Physics, Zhiyuan College.}

\section{Honors and Awards}
\begin{itemize}[itemsep=0pt,topsep=1pt]
  \item National Scholarship for Graduate Students \hfill 2026
  \item Rong Chang Science and Technology Innovation Scholarship \hfill 2024, 2025
  \item Shanghai Outstanding Graduate \hfill 2025
  \item National Scholarship; A-Class Undergraduate Scholarship \hfill 2023
  \item Bronze Medal, Chinese Physics Olympiad (CPhO) \hfill 2019
  \item Jiangsu Provincial First Prize, Chinese Mathematical Olympiad (CMO) \hfill 2019
\end{itemize}

\section{Publications}
{\footnotesize $^*$ Equal contribution. \quad $^\dagger$ Corresponding author.}\par\vspace{4pt}

\publication
  {Quantum Neural Operators with Implicit Quadratic Frame and Expressivity Advantages}
  {\textbf{Ruocheng Wang}$^*$ (co-first author), Xiaoqiu Zhong$^*$, Zhuo Xia, Junchi Yan$^\dagger$.}
  {Nature Machine Intelligence}

\publication
  {CoEvol-NO: State and Coordinate Co-Evolution with an Error-Driven Predictor-Corrector Paradigm for Neural Operator Transformer}
  {Jianqiao Zeng, \textbf{Ruocheng Wang} (second author), Yanzhi Liu, Hao Xiong, Junchi Yan$^\dagger$.}
  {ICML 2026 (Oral)}

\publication
  {QuanONet: Quantum Neural Operator with Application to Differential Equation}
  {\textbf{Ruocheng Wang} (first author), Zhuo Xia, Ge Yan, Junchi Yan$^\dagger$.}
  {ICML 2025}

\publication
  {Rethinking Parity Check Enhanced Symmetry-Preserving Ansatz}
  {Ge Yan, Mengfei Ran, \textbf{Ruocheng Wang}, Kaiseng Pan, Junchi Yan$^\dagger$.}
  {NeurIPS 2024}

\publication
  {Non-Equilibrium and Self-Organization Evolution in Hot-Spot Ignition Processes}
  {X.-Y. Fu$^*$, Z.-Y. Guo$^*$, Q.-H. Wang$^*$, \textbf{R.-C. Wang}$^*$ (co-first author), D. Wu, J. Zhang$^\dagger$.}
  {Physics of Plasmas}

\publication
  {HomeGuard: VLM-based Embodied Safeguard for Identifying Contextual Risk in Household Task}
  {Xiaoya Lu, Yijin Zhou, Zeren Chen, \textbf{Ruocheng Wang}, Bingrui Sima, Enshen Zhou, Lu Sheng, Dongrui Liu, Jing Shao.}
  {ECCV 2026}

\publication
  {AI-Empowered Quantum Circuit Synthesis and Compilation Optimization: Overview and Prospects}
  {Ge Yan, Wenjie Wu, Yuheng Chen, Kaisen Pan, Xudong Lu, Zixiang Zhou, Yuhan Wang, \textbf{Ruocheng Wang}, Junchi Yan$^\dagger$.}
  {TPAMI 2026}

\newpage
\section{Research Projects}

\project{Recurrent RG Net: Renormalization Group Learning without Coupling Labels}
  {2026 -- Present}
\begin{itemize}
  \item Develop a framework for multiscale statistical learning without coupling labels, phase labels, or supervision from known critical points. Jointly learn a local coarse-graining operator $C$, low-dimensional theory coordinates $Q$, and a shared renormalization group map $R$ from configuration ensembles.
  \item Design a joint training framework that combines consistency across ensemble subsets, data geometry across all scales, and one-step and multistep closed-loop constraints to learn finite-scale statistics and their latent-space dynamics.
  \item Completed prototype experiments with one- and two-dimensional latent coordinates on the one-dimensional transverse-field Ising model, qualitatively observing candidate stable, unstable, and saddle fixed points. Ongoing work investigates two-dimensional Ising systems and the reliable identification of critical behavior.
\end{itemize}

\vspace{9pt}
\project{CoEvol-NO: A Neural Operator Transformer with State and Coordinate Co-Evolution}
  {Jan. 2025 -- Jun. 2026}
\begin{itemize}
  \item Proposed an architecture that jointly evolves states and coordinates to address the trade-off between preserving geometric detail and retaining memory of long-term physical evolution in neural operators.
  \item Introduced a classical predictor-corrector paradigm that combines explicit target prediction with error-driven exact gradient correction, enabling the joint evolution of latent variables and mesh sequences.
  \item Proved that the architecture has strictly linear $\mathcal{O}(N)$ time complexity and that its exact gradient correction mechanism overcomes the instability of conventional residual updates under first-order approximation.
  \item Achieved state-of-the-art performance on industrial-scale fluid simulation tasks, including Shape-Net Car and AirfRANS, as well as standard benchmarks such as Navier--Stokes.
\end{itemize}

\vspace{9pt}
\project{QuanONet: Quantum Neural Operators with Theoretical Expressivity Advantages}
  {Dec. 2024 -- Apr. 2026}
\begin{itemize}
  \item Proposed QuanONet, the first quantum neural operator framework for the NISQ era, representing continuous nonlinear mappings as the Frobenius inner product between branch-modulated observables and a trunk density matrix.
  \item Extended the operator universal approximation theorem to the quantum setting and proved that the density matrix implicitly constructs a quadratic feature frame, surpassing the linear capacity limit of classical models to achieve $\mathcal{O}(p^2)$ expressive capacity.
  \item Introduced a trainable-frequency (TF) embedding strategy that adaptively aligns with the dominant eigenfeatures of the output covariance operator, substantially improving the capture of broadband dynamics.
  \item Demonstrated strong optimization robustness in high-dimensional scaling regimes and validated the framework on real IBM superconducting quantum processors.
\end{itemize}

\end{document}
